\section{Diptongos \textit{ui, eu, ei/ij, ou/au}}
\textbf{Qué son:} Combinaciones vocálicas que se pronuncian como un solo núcleo con deslizamiento.

\textbf{Por qué importan:} Son frecuentes y distintivos; \textit{ui} no existe en español.

\textbf{Cómo practicarlos:}
\begin{itemize}[leftmargin=*]
\item \textit{ui}: \textit{huis, muis, luisteren}. Boca en ``o'' breve que se desliza a ``i''.
  \item \textit{eu}: \textit{deur, leuk, keuken}. Similar a \textit{ö} + \textit{u} breve.
  \item \textit{ei/ij}: mismos sonidos, distinto deletreo: \textit{eigen, lijst, klein}. Memoriza por palabra.
  \item \textit{ou/au}: mismo sonido: \textit{houden, vrouw, auto}.
\end{itemize}

\textbf{Ejercicios:}
\begin{itemize}[leftmargin=*]
  \item Parches de repetición: 10 veces cada palabra en ritmo lento-rápido.
  \item Frases cortas: ``\textit{Ik luister naar een leuk huis.}'' ``\textit{Mijn vrouw koopt een auto.}''
\end{itemize}
